% Casper FFG: finalizzazione con un link che salta un'epoca.
% C_N e C_{N+1} sono giustificati; il link C_N -> C_{N+2} con più dei 2/3 dei voti
% giustifica C_{N+2} e finalizza C_N.
% Incluso nel documento con % Casper FFG: finalizzazione con un link che salta un'epoca.
% C_N e C_{N+1} sono giustificati; il link C_N -> C_{N+2} con più dei 2/3 dei voti
% giustifica C_{N+2} e finalizza C_N.
% Incluso nel documento con % Casper FFG: finalizzazione con un link che salta un'epoca.
% C_N e C_{N+1} sono giustificati; il link C_N -> C_{N+2} con più dei 2/3 dei voti
% giustifica C_{N+2} e finalizza C_N.
% Incluso nel documento con % Casper FFG: finalizzazione con un link che salta un'epoca.
% C_N e C_{N+1} sono giustificati; il link C_N -> C_{N+2} con più dei 2/3 dei voti
% giustifica C_{N+2} e finalizza C_N.
% Incluso nel documento con \input{immagini/finalizzazione-salto.tex}.
\definecolor{fsfinal}{HTML}{2E9E5B}
\definecolor{fsgiust}{HTML}{5B9BD5}
\definecolor{fsepoca}{HTML}{F2F2F2}
\definecolor{fsvoto}{HTML}{C55A11}
\begin{tikzpicture}[
    slot/.style={draw, fill=white, minimum size=0.42cm, inner sep=0pt},
    checkpoint/.style={draw, minimum size=0.75cm, inner sep=0pt,
        font=\small\bfseries, text=white},
    puntatore/.style={-{Latex[length=1.6mm]}, line width=0.6pt},
    voto/.style={-{Latex[length=2.5mm]}, line width=1.1pt, dashed, fsvoto}
]
    \foreach \k/\nome/\colore/\stato in {0/N/fsfinal/finalizzato,
            1/N+1/fsgiust/giustificato, 2/N+2/fsgiust/giustificato} {
        \begin{scope}[shift={(\k*5.6, 0)}]
            \draw[fill=fsepoca, rounded corners=3pt] (0, -0.5) rectangle (5.1, 0.5);
            \node[font=\small] at (3.1, -0.8) {epoca $\nome$};
            \node[checkpoint, fill=\colore] (c\k) at (0.6, 0) {$C_{\nome}$};
            \node[slot] (a\k) at (1.6, 0) {};
            \node[slot] (b\k) at (2.4, 0) {};
            \node[slot] (d\k) at (3.2, 0) {};
            \node (p\k) at (3.9, 0) {\dots};
            \node[slot] (z\k) at (4.6, 0) {};
            \node[font=\scriptsize, text=\colore, below=2pt] at (c\k.south |- 0, -0.5) {\stato};
            \draw[puntatore] (a\k) -- (c\k);
            \draw[puntatore] (b\k) -- (a\k);
            \draw[puntatore] (d\k) -- (b\k);
            \draw[puntatore] (p\k) -- (d\k);
            \draw[puntatore] (z\k) -- (p\k);
        \end{scope}
    }
    \draw[puntatore] (c1) -- (z0);
    \draw[puntatore] (c2) -- (z1);

    % il link salta C_{N+1}
    \draw[voto] (c0.north) -- ++(0, 0.9) -| node[pos=0.25, above, font=\small]
        {voti $> 2/3$ dello stake} (c2.north);
\end{tikzpicture}
.
\definecolor{fsfinal}{HTML}{2E9E5B}
\definecolor{fsgiust}{HTML}{5B9BD5}
\definecolor{fsepoca}{HTML}{F2F2F2}
\definecolor{fsvoto}{HTML}{C55A11}
\begin{tikzpicture}[
    slot/.style={draw, fill=white, minimum size=0.42cm, inner sep=0pt},
    checkpoint/.style={draw, minimum size=0.75cm, inner sep=0pt,
        font=\small\bfseries, text=white},
    puntatore/.style={-{Latex[length=1.6mm]}, line width=0.6pt},
    voto/.style={-{Latex[length=2.5mm]}, line width=1.1pt, dashed, fsvoto}
]
    \foreach \k/\nome/\colore/\stato in {0/N/fsfinal/finalizzato,
            1/N+1/fsgiust/giustificato, 2/N+2/fsgiust/giustificato} {
        \begin{scope}[shift={(\k*5.6, 0)}]
            \draw[fill=fsepoca, rounded corners=3pt] (0, -0.5) rectangle (5.1, 0.5);
            \node[font=\small] at (3.1, -0.8) {epoca $\nome$};
            \node[checkpoint, fill=\colore] (c\k) at (0.6, 0) {$C_{\nome}$};
            \node[slot] (a\k) at (1.6, 0) {};
            \node[slot] (b\k) at (2.4, 0) {};
            \node[slot] (d\k) at (3.2, 0) {};
            \node (p\k) at (3.9, 0) {\dots};
            \node[slot] (z\k) at (4.6, 0) {};
            \node[font=\scriptsize, text=\colore, below=2pt] at (c\k.south |- 0, -0.5) {\stato};
            \draw[puntatore] (a\k) -- (c\k);
            \draw[puntatore] (b\k) -- (a\k);
            \draw[puntatore] (d\k) -- (b\k);
            \draw[puntatore] (p\k) -- (d\k);
            \draw[puntatore] (z\k) -- (p\k);
        \end{scope}
    }
    \draw[puntatore] (c1) -- (z0);
    \draw[puntatore] (c2) -- (z1);

    % il link salta C_{N+1}
    \draw[voto] (c0.north) -- ++(0, 0.9) -| node[pos=0.25, above, font=\small]
        {voti $> 2/3$ dello stake} (c2.north);
\end{tikzpicture}
.
\definecolor{fsfinal}{HTML}{2E9E5B}
\definecolor{fsgiust}{HTML}{5B9BD5}
\definecolor{fsepoca}{HTML}{F2F2F2}
\definecolor{fsvoto}{HTML}{C55A11}
\begin{tikzpicture}[
    slot/.style={draw, fill=white, minimum size=0.42cm, inner sep=0pt},
    checkpoint/.style={draw, minimum size=0.75cm, inner sep=0pt,
        font=\small\bfseries, text=white},
    puntatore/.style={-{Latex[length=1.6mm]}, line width=0.6pt},
    voto/.style={-{Latex[length=2.5mm]}, line width=1.1pt, dashed, fsvoto}
]
    \foreach \k/\nome/\colore/\stato in {0/N/fsfinal/finalizzato,
            1/N+1/fsgiust/giustificato, 2/N+2/fsgiust/giustificato} {
        \begin{scope}[shift={(\k*5.6, 0)}]
            \draw[fill=fsepoca, rounded corners=3pt] (0, -0.5) rectangle (5.1, 0.5);
            \node[font=\small] at (3.1, -0.8) {epoca $\nome$};
            \node[checkpoint, fill=\colore] (c\k) at (0.6, 0) {$C_{\nome}$};
            \node[slot] (a\k) at (1.6, 0) {};
            \node[slot] (b\k) at (2.4, 0) {};
            \node[slot] (d\k) at (3.2, 0) {};
            \node (p\k) at (3.9, 0) {\dots};
            \node[slot] (z\k) at (4.6, 0) {};
            \node[font=\scriptsize, text=\colore, below=2pt] at (c\k.south |- 0, -0.5) {\stato};
            \draw[puntatore] (a\k) -- (c\k);
            \draw[puntatore] (b\k) -- (a\k);
            \draw[puntatore] (d\k) -- (b\k);
            \draw[puntatore] (p\k) -- (d\k);
            \draw[puntatore] (z\k) -- (p\k);
        \end{scope}
    }
    \draw[puntatore] (c1) -- (z0);
    \draw[puntatore] (c2) -- (z1);

    % il link salta C_{N+1}
    \draw[voto] (c0.north) -- ++(0, 0.9) -| node[pos=0.25, above, font=\small]
        {voti $> 2/3$ dello stake} (c2.north);
\end{tikzpicture}
.
\definecolor{fsfinal}{HTML}{2E9E5B}
\definecolor{fsgiust}{HTML}{5B9BD5}
\definecolor{fsepoca}{HTML}{F2F2F2}
\definecolor{fsvoto}{HTML}{C55A11}
\begin{tikzpicture}[
    slot/.style={draw, fill=white, minimum size=0.42cm, inner sep=0pt},
    checkpoint/.style={draw, minimum size=0.75cm, inner sep=0pt,
        font=\small\bfseries, text=white},
    puntatore/.style={-{Latex[length=1.6mm]}, line width=0.6pt},
    voto/.style={-{Latex[length=2.5mm]}, line width=1.1pt, dashed, fsvoto}
]
    \foreach \k/\nome/\colore/\stato in {0/N/fsfinal/finalizzato,
            1/N+1/fsgiust/giustificato, 2/N+2/fsgiust/giustificato} {
        \begin{scope}[shift={(\k*5.6, 0)}]
            \draw[fill=fsepoca, rounded corners=3pt] (0, -0.5) rectangle (5.1, 0.5);
            \node[font=\small] at (3.1, -0.8) {epoca $\nome$};
            \node[checkpoint, fill=\colore] (c\k) at (0.6, 0) {$C_{\nome}$};
            \node[slot] (a\k) at (1.6, 0) {};
            \node[slot] (b\k) at (2.4, 0) {};
            \node[slot] (d\k) at (3.2, 0) {};
            \node (p\k) at (3.9, 0) {\dots};
            \node[slot] (z\k) at (4.6, 0) {};
            \node[font=\scriptsize, text=\colore, below=2pt] at (c\k.south |- 0, -0.5) {\stato};
            \draw[puntatore] (a\k) -- (c\k);
            \draw[puntatore] (b\k) -- (a\k);
            \draw[puntatore] (d\k) -- (b\k);
            \draw[puntatore] (p\k) -- (d\k);
            \draw[puntatore] (z\k) -- (p\k);
        \end{scope}
    }
    \draw[puntatore] (c1) -- (z0);
    \draw[puntatore] (c2) -- (z1);

    % il link salta C_{N+1}
    \draw[voto] (c0.north) -- ++(0, 0.9) -| node[pos=0.25, above, font=\small]
        {voti $> 2/3$ dello stake} (c2.north);
\end{tikzpicture}
